\section{Giàn giáo, tự giải thích và lỗi quan niệm}

\subsection{Giàn giáo và nghịch lý hỗ trợ}
Giàn giáo là sự hỗ trợ tạm thời giúp người học làm được điều mình chưa tự làm được, và sự hỗ trợ đó được rút dần khi người học làm chủ nhiệm vụ~\cite{wood1976}. Khái niệm này gắn với vùng phát triển gần của Vygotsky, tức khoảng giữa điều học sinh tự làm được và điều học sinh làm được khi có người giúp~\cite{vygotsky1978}. Hình~\ref{fig:zpd} minh họa mối quan hệ đó. Gia sư hiệu quả luôn nhắm vào vùng giữa: nhiệm vụ nằm trong vùng tự làm được thì không cần giúp, còn nhiệm vụ nằm ngoài vùng phát triển gần thì sự giúp đỡ không đủ để học sinh tiến bộ.

\begin{figure}[ht]
\centering
\begin{tikzpicture}[font=\footnotesize]
\fill[gray!8] (0,0) ellipse (5.0cm and 3.1cm);
\fill[gray!22] (0,0) ellipse (3.4cm and 2.2cm);
\fill[gray!42] (0,0) ellipse (1.6cm and 0.8cm);
\draw[line width=0.6pt] (0,0) ellipse (5.0cm and 3.1cm);
\draw[line width=0.6pt] (0,0) ellipse (3.4cm and 2.2cm);
\draw[line width=0.6pt] (0,0) ellipse (1.6cm and 0.8cm);
\node[align=center] at (0,0) {\textbf{Tự làm được}};
\node[align=center] at (0,1.5) {\textbf{Làm được khi có hỗ trợ}\\\scriptsize\itshape vùng phát triển gần};
\node[align=center] at (0,-2.65) {\textbf{Chưa làm được}};
\draw[-{Stealth[length=2mm,width=1.6mm]},line width=0.7pt] (-7.3,0.9) -- (-1.9,0.3);
\node[font=\scriptsize\itshape,anchor=south west] at (-7.5,1.0) {giàn giáo hỗ trợ};
\draw[-{Stealth[length=2mm,width=1.6mm]},line width=0.7pt] (7.3,-0.9) -- (3.5,-0.9);
\node[font=\scriptsize\itshape,anchor=north east] at (7.5,-1.0) {rút dần hỗ trợ};
\end{tikzpicture}

\caption{Vùng phát triển gần. Giàn giáo đưa nhiệm vụ từ vùng làm được khi có hỗ trợ vào vùng tự làm được, và hỗ trợ được rút dần khi năng lực tăng.}
\label{fig:zpd}
\end{figure}

Koedinger và Aleven gọi mâu thuẫn trong việc quyết định mức hỗ trợ là nghịch lý hỗ trợ: giúp quá nhiều làm học sinh mất cơ hội tự nghĩ, giúp quá ít làm học sinh bế tắc~\cite{koedinger2007}. Nghiên cứu về hành vi xin trợ giúp cho thấy học sinh thường lạm dụng gợi ý hoặc né tránh gợi ý, nên hệ thống cần quy tắc rõ ràng về thời điểm và mức độ hỗ trợ~\cite{aleven2016}. Các hệ gia sư thường dùng thang gợi ý đi từ chung đến cụ thể, với bậc cuối là chính lời giải. Đồ án dùng một thang tám bậc từ H0 đến H7 theo hai trục là bậc (giúp bao nhiêu) và mục tiêu (giúp vào đâu). Mỗi bậc có một trọng số độc lập do khung đánh giá quy định. Gia sư và bộ chấm dùng chung các trọng số đó, nên hai bên cùng một thước đo về mức học sinh đã được giúp. Chương~\ref{chap:design} mô tả chi tiết thang này.

\subsection{Tự giải thích và hiệu ứng đảo ngược chuyên môn}
Hiệu ứng tự giải thích cho thấy học sinh tự giải thích các bước và lý do của mình thì hiểu sâu và nhớ lâu hơn~\cite{chi1989,chi1994}. Chi và cộng sự phân tích việc học từ gia sư như quá trình học sinh tự xây dựng kiến thức chứ không phải tiếp nhận thụ động~\cite{chi2001}. Kết quả này là cơ sở của hai nước đi trong đồ án, một nước đi hỏi ``vì sao'' và một nước đi mời học sinh giảng lại cách làm.

Hiệu ứng đảo ngược chuyên môn đặt giới hạn cho việc áp dụng ở trên. Một kiểu hỗ trợ có ích với người mới có thể không giúp được, thậm chí gây cản trở, với người đã thành thạo~\cite{kalyuga2003}. Gia sư vì vậy chỉ hỏi ``vì sao'' khi kiến thức còn mới hoặc chưa vững, và hỏi ít đi khi học sinh đã vững.

\subsection{Lỗi quan niệm và xung đột nhận thức}
Lỗi quan niệm là một cách hiểu sai có hệ thống chứ không phải sai ngẫu nhiên. Resnick và cộng sự ghi nhận ở số thập phân các lỗi điển hình như cho rằng số nào dài hơn thì lớn hơn (viết $0{,}25>0{,}3$), do học sinh chuyển cách nghĩ của số tự nhiên sang~\cite{resnick1989}. Một lỗi quen thuộc khác là nhầm hàng phần mười với hàng phần trăm, chẳng hạn viết $9/10$ thành $0{,}09$.

Xung đột nhận thức bắt nguồn từ lý thuyết thay đổi khái niệm: người dạy đưa cho học sinh một phản ví dụ để học sinh tự thấy cách nghĩ của mình dẫn tới mâu thuẫn, thay vì nói thẳng rằng cách nghĩ đó sai~\cite{posner1982}. Đồ án gắn mỗi lỗi quan niệm mà giáo viên khai báo với một phản ví dụ và dùng phản ví dụ đó như một nước đi sư phạm riêng. Bộ nhận diện lỗi quan niệm là một luật tất định dựa trên dạng sai đặc trưng, không dựa trên phán đoán của mô hình ngôn ngữ.

\subsection{Khen nỗ lực}
Mueller và Dweck cho thấy khen trí thông minh làm trẻ ngại thử thách và nản khi gặp khó, còn khen nỗ lực thì cho kết quả ngược lại~\cite{mueller1998}. Bộ kiểm chứng của đồ án vì thế chặn lời khen trí thông minh. Bộ kiểm chứng cũng chặn lời khen đứng trước một câu trả lời sai, vì lời khen sai chỗ làm học sinh hiểu nhầm rằng câu trả lời đã đúng.
